% =====================================================================
\section{Conclusion}\label{sec:conclusion}
% =====================================================================
This paper presented GarbageGo, a domain AI agent that changes the
question residents ask from ``when will the truck arrive'' to ``which
point can I still make it to now.'' The scheme combines three elements:
a slack measure that compares truck ETA with network walking time and a
buffer that grows with data age and remaining stops; a fallback to
clearly labeled schedules when live data is stale; and a division of
labor in which the LLM explains results but never produces a number.
A 29-day study of Taichung's open data shows that the required inputs
are available: license plates link 96.1\% of collection points to
vehicles, and local geocoding located 95.6\% of point addresses.

The same study also defines the limits of this work. First, the vehicle
feed delivers snapshots only about every five minutes, and 18\% of
vehicles in one snapshot had been silent for over 30 minutes, so ETA
inputs can be stale. Second, no ground-truth arrival times exist yet, so
no recommendation accuracy is claimed. Third, the buffer coefficients in
(\ref{eq:buffer}) are not yet calibrated, and 4.4\% of point addresses
still lack verified coordinates.

Future work follows directly from these limits. We will validate how
actual arrival times can be inferred from five-minute snapshots,
calibrate the buffer on the ongoing data collection, and compare the ETA
baselines by MAE and RMSE. Finally, we will measure whether slack-based
recommendation raises the catch rate and lowers the misjudgment rate
compared with the nearest-point strategy.
